\section{Выводы}

В данной статье было проведено сравнение двух моделей (LSA и MiniLM)
по метрикам STS~$\rho$, Accuracy, ARI, PR, ER, IS и HS.
Сравнение производилось на датасетах STSb, 20~Newsgroups и MRPC.
Результаты показывают различия геометрических свойств пространств
и качества обеих моделей на практических задачах: более
сфокусированные главные направления и превосходство в семантической
обработке у MiniLM и равномерно распределённые направления без учёта
контекста у LSA, что ограничивает её в обработке контекстно-зависимых
явлений, таких как полисемия~\cite{hofmann1999plsa}.
Двунаправленное контекстное предобучение MiniLM позволяет преодолеть
это ограничение~\cite{artetxe2022bidirectionality}.
В задачах классификации и нахождения парафраз LSA
сопоставима с MiniLM, так как эти задачи имеют ключевые слова и не
требуют контекста. Вместе с тем сравнение имеет ограничения: MiniLM
предобучен на большом корпусе, размерности моделей различаются, а
каждый класс представлен одной моделью.
